% =============================================================================
% Chapter 6 --- Simulation Results
% Every number in this chapter is a direct quote from a file under target/.
% =============================================================================

\chapter{Simulation Results}
\label{ch:results}

This chapter consolidates the simulation evidence. Every number quoted in
the text or in a table is a direct read of a report file under
\sv{target/}. \Cref{tab:results-summary} gives the one-line verdict for
each testbench; subsequent sections walk through the evidence in more
detail.

\begin{table}[htbp]
    \centering
    \caption{Summary of checked-in simulation outcomes. Counts and
        summary strings are read verbatim from the corresponding report
        files.}
    \label{tab:results-summary}
    \begin{tabularx}{\linewidth}{l l X}
        \toprule
        \textbf{Testbench} & \textbf{Outcome} & \textbf{Evidence} \\
        \midrule
        \sv{controller\_addr\_pulse\_tb} & 30 tests, all PASS &
            30 ``\sv{PASS:}'' lines in
            \sv{target/Controller/controller\_addr\_pulse\_verify.txt}. \\
        \sv{parallel\_interface\_controller\_integration\_tb} &
            \sv{PASS=10 FAIL=0} &
            Summary in
            \sv{target/Controller/tb\_pi\_controller\_integration.txt}. \\
        \sv{controller\_prog\_verify\_lut\_tb} &
            0 errors in 640 weights &
            Summary in
            \sv{target/Controller/prog/prog\_verify\_report.txt};
            additional asymmetric-flow coverage line. \\
        \sv{controller\_prog\_verify\_lut\_10w\_tb} &
            10 PASS, 0 errors &
            Summary in
            \sv{target/Controller/prog/prog\_verify\_10w\_report.txt}. \\
        \sv{controller\_host\_erase\_tb} &
            PASS, cell A cleared, cell B preserved &
            \sv{target/Controller/erase/controller\_host\_erase\_report.txt}. \\
        \sv{controller\_host\_read\_reorder\_tb} &
            \sv{PASS=8 FAIL=0} &
            \sv{target/Controller/read/controller\_host\_read\_reorder\_report.txt}. \\
        \sv{controller\_inf\_buffer\_flow\_tb} & PASS &
            \sv{target/Controller/inf/controller\_inf\_buffer\_flow\_tb\_log.txt}. \\
        Input-buffer unit benches &
            Pass-oriented summaries &
            Logs in \sv{target/Input\_Buffer/}. \\
        Parallel-interface unit bench & 6 checks pass &
            \sv{target/Parallel\_Interface/parallel\_interface\_extract\_tb\_log.txt}. \\
        \bottomrule
    \end{tabularx}
\end{table}

\section{Address and Pulse-Mode Packing (30 Directed Cases)}
\label{sec:results-addrpulse}

\sv{controller\_addr\_pulse\_tb.sv} runs 30 directed tests covering
\sv{CMD\_PROG}, \sv{CMD\_READ}, and \sv{CMD\_ERASE} across the address
space, with randomized PE, SA, column, and row indices and randomized
data bytes. For each test, the bench computes the expected outgoing
\signame{ann\_core\_word} and the expected \signame{pulses} encoding
from \sv{controller\_pkg} helpers, and compares them to the live RTL
outputs. The report header (\sv{controller\_addr\_pulse\_verify.txt})
documents the parameters in force:

\begin{lstlisting}[style=plainstyle,language={}]
// Controller Address and Pulse Verification
// Parameters: TREAD=2, PULSE_NUM_READ=1, TPROG=2, PULSE_NUM_PROG=1
//            TERASE=2, PULSE_NUM_ERASE=1, TINF=2, PULSE_NUM_INF=1
// Expected: READ->pulses=001, PROG->pulses=010,
//           ERASE->pulses=011, INF->pulses=100
\end{lstlisting}

All 30 tests print a \sv{PASS:} line with the expected word and its
binary field decomposition, for example

\begin{lstlisting}[style=plainstyle,language={}]
=== Test 2: READ   (PE=1,SA=3,col=6,row=7) at addr 0x01e7 ===
  PASS: ann_core_word 0x00281080
         (0b00000000-0010-1000-00010000-10000000)
\end{lstlisting}

\noindent The test confirms that the left-shift one-hot reconstruction
inside \sv{pack\_ann\_core\_word} agrees with the package helpers that
generate the stimulus, and that \signame{pulses} takes the expected
\sv{PULSE\_MODE\_*} value for each command class. This is the static
packing contract of the controller.

\section{Parallel-Interface / Controller / Input-Buffer Integration (10 Cases)}
\label{sec:results-integration}

\sv{parallel\_interface\_controller\_integration\_tb.sv} drives ten
labelled scenarios (\sv{01\_READ} through \sv{10\_ERASE}), mixing
\sv{CMD\_READ}, \sv{CMD\_PROG}, \sv{CMD\_ERASE}, and \sv{CMD\_INF}. For
each case it records, in the corresponding block of
\sv{tb\_pi\_controller\_integration.txt}, the host packet, the
PI-decoded bus, the DUT busy duration, the expected and observed pulse
mode, the expected and observed \signame{ann\_core\_word}, and an 8-by-8
snapshot of the input buffer. The file ends with

\begin{lstlisting}[style=plainstyle,language={}]
// Summary: PASS=10 FAIL=0
\end{lstlisting}

A representative case (case \sv{02\_PROG}) shows a host packet with
\signame{host\_data}~=~\sv{0x0c110101} and \signame{host\_cmd}~=~\sv{CMD\_PROG}
being accepted at the interface (decoded to
\signame{valid}~=~1, \signame{data}~=~\sv{0x0c},
\signame{address}~=~\sv{0x0000}), the DUT staying busy for 22 posedge
cycles while exercising PROG and VERIFY, and the input-buffer snapshot
showing \sv{0c} deposited at the zero address ---
reflecting the \sv{CTRL\_DATA\_LOAD} write during \sv{PROG\_HIZ}/\sv{PROG\_SELECT}
described in \Cref{sec:rtl-controller-mainfsm}. Case \sv{04\_INF} shows
the inference burst: the row~0 of the buffer ends up filled with eight
copies of the payload byte \sv{0xaa}, because the nine-beat
\sv{CMD\_INF} sequence writes eight consecutive pixels and then
transitions to \stateval{S\_COMPUTE}.

\section{640-Weight Program--Verify Regression}
\label{sec:results-640}

\sv{controller\_prog\_verify\_lut\_tb.sv} is the primary regression for
the recovery policy described in \Cref{sec:rtl-controller-verify-check}.
It loads 640 expected weights from
\sv{target/Controller/programming\_inputs/weight\_matrix.txt}, programs
every cell under the LUT-based first-attempt schedule, and applies the
fault-injection gating described in \Cref{sec:verif-injection}. The
report \sv{prog\_verify\_report.txt} ends with

\begin{lstlisting}[style=plainstyle,language={}]
// Total weights programmed: 640

// Asymmetric flow coverage:
//   ERASE path (read>expected): 14
//   re-PROG path (read<expected): 3017 (PROG after VERIFY)
// Summary: 0 errors in 640 weights
\end{lstlisting}

The two coverage numbers quantify what the fault-injection machinery
achieved: 14 transitions into \stateval{S\_ERASE} via the verify-failure
branch (over-programmed readbacks) and 3017 re-programming events after
verify (under-programmed readbacks, each resolved within the retry
budget or leading to erase-and-retry per the FSM). The
\sv{0 errors in 640 weights} statement is the end-of-run scoreboard
comparison between the shadow matrix and the expected nibbles loaded
from \sv{weight\_matrix.txt}.

The phase log also contains readable segments such as the following
(clipped from the report):

\begin{lstlisting}[style=plainstyle,language={}]
[314000] Phase: PROG  ann_core_word=0x0a110801 programmed_weight=10
// re-PROG sequence START [326000]  (after phase: CHECK_DONE)
[326000] Phase: REPROG ann_core_word=0x0a110801 programmed_weight=10
[338000] Phase: REPROG ann_core_word=0x0a110801 programmed_weight=10
[350000] Phase: REPROG ann_core_word=0x0a110801 programmed_weight=10
...
[450000] Phase: VERIFY ann_core_word=0x0a110801
         expected_weight=10  read_weight=10
// re-PROG sequence END [458000]  (DUT idle)
\end{lstlisting}

This trace shows an under-programmed injection at weight index 4 (value
10): an initial \sv{PROG} phase is followed by a sequence of \sv{REPROG}
phases under the short retry train described in
\Cref{sec:rtl-controller-pulses}, culminating in a \sv{VERIFY} that
matches the expected value and returns the DUT to idle.

\section{Compact 10-Weight Program--Verify}
\label{sec:results-10w}

\sv{controller\_prog\_verify\_lut\_10w\_tb.sv} is a compact variant used
for waveform-driven debugging. Its header enumerates the scenarios per
index:

\begin{lstlisting}[style=plainstyle,language={}]
// idx0 equal | idx1 read< | idx2 read> |
// idx3-4,6-9 happy | idx5 max re-PROG (prog_retry_cnt) -> ERASE
\end{lstlisting}

and the report ends with

\begin{lstlisting}[style=plainstyle,language={}]
  PASS idx=0 addr=0x0000 weight=2
  PASS idx=1 addr=0x0008 weight=2
  PASS idx=2 addr=0x0010 weight=5
  PASS idx=3 addr=0x0018 weight=10
  PASS idx=4 addr=0x0020 weight=12
  PASS idx=5 addr=0x0028 weight=10
  PASS idx=6 addr=0x0030 weight=5
  PASS idx=7 addr=0x0038 weight=2
  PASS idx=8 addr=0x0040 weight=2
  PASS idx=9 addr=0x0048 weight=2
// Summary: 0 errors
\end{lstlisting}

so all four per-index scenarios (match, under, over, max-retry-to-erase)
land on the expected value. This bench is the one the user guide
recommends for fast iteration on FSM or pulse-shaping changes; its
short runtime makes it practical to open in the ModelSim GUI with the
companion wave \sv{.do} scripts.

\section{Host Erase and Non-Target Retention}
\label{sec:results-erase}

\sv{controller\_host\_erase\_tb.sv} programs two cells at distinct
addresses, then issues a single host \sv{CMD\_ERASE} on one of them.
The bench's pulse trace records the DUT's state transitions; the
\sv{state} column shows \sv{2} (\stateval{S\_PROGRAM}) and \sv{3}
(\stateval{S\_VERIFY}) during the PROG phases, then \sv{4}
(\stateval{S\_ERASE}) during the host-erase phase. The erase phase
advances through its sub-states (\sv{erase\_state} values
\sv{0}\,\(\rightarrow\)\,\sv{1}\,\(\rightarrow\)\,\sv{2}\,\(\rightarrow\)\,\sv{3}\,\(\rightarrow\)\,\sv{4}
in the log), with \signame{pulses}~=~\sv{3'b011} (\sv{PULSE\_MODE\_ERASE})
asserted during \sv{ERASE\_PULSE}, and returns to idle without entering
\stateval{S\_PROGRAM} --- consistent with the \signame{erase\_from\_host}
branch of \Cref{tab:erase-complete}. The report concludes with:

\begin{lstlisting}[style=plainstyle,language={}]
// AFTER ERASE (expect A=0, B unchanged)
// Cell A: addr=0x01e7  indices [1][3][7][4]  weight=0
// Cell B: addr=0x0100  indices [1][0][0][0]  weight=4

PASS: cell A cleared, cell B still 4
\end{lstlisting}

demonstrating, within the limits of the behavioral mock, that a host
erase affects only the addressed cell.

\section{Host Read Reorder (8 Permuted Reads)}
\label{sec:results-readreorder}

\sv{controller\_host\_read\_reorder\_tb.sv} programs eight distinct
cells with distinct weights, then issues \sv{CMD\_READ} in the
permuted program-index order \sv{4 0 7 2 1 6 3 5}. For each read the
bench samples \signame{ann\_core\_word} while the DUT is in
\stateval{S\_READ}, reads back the mocked \signame{weight\_read\_data},
and emits a PASS/FAIL line. Every one of the eight reads passes, and the
report ends with

\begin{lstlisting}[style=plainstyle,language={}]
// Summary: PASS=8 FAIL=0
\end{lstlisting}

The out-of-order scenario is relevant because it exercises
\signame{address\_reg} capture on every new transaction: stale
registered address fields from the previous program cycle would produce
the wrong one-hot tail on \signame{ann\_core\_word} and would be caught
by the comparison.

\section{Inference Flow}
\label{sec:results-inf}

\sv{controller\_inf\_buffer\_flow\_tb.sv} drives the INF path through
\stateval{S\_COLLECT\_DATA}\,\(\rightarrow\)\,\stateval{S\_COMPUTE}.
Across the collect phase, the bench increments \signame{writes\_seen}
whenever \signame{buf\_reg\_ctrl}~=~\sv{CTRL\_DATA\_LOAD} and
\signame{buf\_read\_write}~=~1; it explicitly fails if a compute
indicator (\sv{CTRL\_COMPUTE} or \signame{pulses}~=~\sv{PULSE\_MODE\_INF})
fires before eight writes have been seen. Once eight pixels are in the
buffer, an optional extra INF beat drives the FSM into
\stateval{S\_COMPUTE}, where the bench samples the bit-serial outputs
\signame{D0}\ldots\signame{D7} against expected column bits derived
from a row snapshot; each matching lane is counted as a pass and each
mismatch as a fail in the per-lane checker. The log in
\sv{target/Controller/inf/controller\_inf\_buffer\_flow\_tb\_log.txt}
shows a pass-oriented summary.

\section{Unit Benches}
\label{sec:results-unit}

The three input-buffer benches exercise bit-serial output correctness,
reset-to-zero behavior, and full-capacity overwrite behavior. Each
writes a \sv{SUMMARY: N passed, M failed} line into its log under
\sv{target/Input\_Buffer/}, reflecting the directed check counts
in the bench body. The parallel-interface bench
\sv{parallel\_interface\_extract\_tb.sv} executes six directed checks
of valid/invalid tail decode and address extraction; its log under
\sv{target/Parallel\_Interface/} carries a similar summary line.

\section{Scope of the Reported Evidence}
\label{sec:results-scope}

Before moving on to the interpretive chapter, it is useful to be
explicit about what \Cref{tab:results-summary} and the numbers quoted
in this chapter actually prove:

\begin{itemize}
    \item Under the behavioral memristor model and the
          \signame{op\_done} generator described in \Cref{ch:verification},
          the controller \emph{does} program 640 cells with zero
          end-of-run mismatches against the expected weight matrix,
          while exercising the recovery branches of the FSM in both
          directions.
    \item The host-to-core static contract (address packing, pulse-mode
          decoding) is verified by directed unit- and integration-level
          checks.
    \item The inference path correctly stages pixel writes before
          triggering compute, and the bit-serial output lanes match the
          expected per-bit values drawn from the buffer snapshot.
\end{itemize}

\noindent And, equally important, what they do not prove:

\begin{itemize}
    \item Timing closure or physical feasibility --- no synthesis or
          place-and-route is in scope.
    \item End-to-end numerical correctness of an MVM against an
          analytical reference --- no golden MAC is in the environment.
    \item Behavior under analog non-idealities not captured by a 4-bit
          integer shadow matrix.
\end{itemize}

These scope limits feed directly into \Cref{ch:analysis} (what the
results prove, interpreted), and \Cref{ch:limitations} (what remains
to be done).
